\subsection{Objective}

Task 4 requires converting the Task~3 sub-image of size $899\times148$ into a
binary image and showing the intermediate results.  Three thresholding schemes
were compared, and the selected method keeps the ten characters as separate
foreground components while rejecting most of the textured background.

\subsection{Binarization Methods Compared}

\begin{enumerate}
  \item \textbf{Global Otsu thresholding.}  Otsu's method~\cite{otsu1979} chooses the level
  $T^{*}$ that maximizes the between-class variance of the intensity
  histogram:
  \begin{equation}
    T^{*} = \arg\max_{T}\,
    w_{0}(T)\,w_{1}(T)\bigl(\mu_{0}(T)-\mu_{1}(T)\bigr)^{2},
    \qquad B(x,y) = I(x,y) > T^{*},
    \label{eq:task4-otsu}
  \end{equation}
  where $w_{k}$ and $\mu_{k}$ are the weight and mean of each class.  A single
  global level is simple, but it cannot adapt to the spatially varying
  background texture.

  \item \textbf{Global Otsu after smoothing and morphological opening.}  The
  sub-image is first smoothed with the $5\times5$ box filter from Task~2, then
  thresholded globally, and finally opened with a $3\times3$ structuring
  element (an erosion followed by a dilation) to delete small speckle
  components.

  \item \textbf{Local adaptive mean thresholding.}  The threshold becomes a
  slowly varying field based on the local mean over a $41\times41$ window:
  \begin{equation}
    T(x,y) = \mu_{41}(x,y) + 25,
    \qquad B(x,y) = I(x,y) > T(x,y).
    \label{eq:task4-adaptive}
  \end{equation}
  After $5\times5$ smoothing, a pixel is foreground only when it is more than
  25 grey levels brighter than its local surroundings.  This is well matched to
  bright characters on a dark, locally textured background.
\end{enumerate}

\begin{figure}[!htbp]
  \centering
  \includegraphics[width=\textwidth]{task4_comparison.png}
  \caption{The Task~3 sub-image and the three binarization candidates.  White
  pixels are foreground.}
  \label{fig:task4-comparison}
\end{figure}

\begin{table}[!htbp]
  \centering
  \caption{Task 4 binarization comparison.}
  \label{tab:task4-metrics}
  \begin{tabular}{>{\raggedright\arraybackslash}p{3.5cm}cccc}
    \toprule
    Method & Foreground (\%) & Components & Character comp. & Max.\ noise (px) \\
    \midrule
    Global Otsu, $T=125$ & 34.00 & 296 & 7 & 131 \\
    Smoothed Otsu $+$ open $3\times3$ & 34.31 & 32 & 7 & 135 \\
    Adaptive $41\times41$, $+25$ (chosen) & 30.72 & 126 & \textbf{10} & 224 \\
    \bottomrule
  \end{tabular}
\end{table}

\subsection{Selected Method and Analysis}

Global Otsu thresholding at $T=125$ performs poorly here.  It labels 34.0\% of
the image as foreground because the bright ring texture crosses the single
global level, producing 296 connected components, and its character strokes
merge into only 7 character-sized components.  Smoothing followed by opening
reduces the component count to 32, but the characters remain merged and the
largest speckles survive.

The adaptive mean threshold is selected as the final method.  It labels 30.7\%
of the image as foreground and produces 126 components, of which exactly 10
have an area above 500 pixels; these ten components correspond one-to-one to
the ten characters of \texttt{HD44780A00} as read left to right.  The remaining
116 components are noise smaller than 225 pixels, which are easy to remove with
an area criterion in Task~6.  Unlike Otsu, the local criterion does not require
one grey level to separate character from texture across the whole image;
instead it marks only pixels that stand out against their own neighbourhood.

The selected binary image is saved as \texttt{task\_4/outputs/task4\_binary.png}
and becomes the input for the outline detection and segmentation tasks.

\FloatBarrier
